\chapter{Proposta}
\label{cap:proposta}

% wiki: Thesis Goal, Pipeline Design, LLM Stage Design, Soundness Assumptions,
% Dashboard Design; project-src/contracts/README.md

Este capítulo descreve a ferramenta proposta. A Seção~\ref{sec:visao-geral}
apresenta o pipeline completo e a Seção~\ref{sec:principio}, a regra que
orienta todas as decisões de projeto. As seções seguintes detalham a entrada, a
etapa estática, o registro de contexto, a etapa com \llm, a re-verificação, os
contratos entre as partes e as premissas sob as quais a ferramenta não perde
defeitos.

\section{Visão geral}
\label{sec:visao-geral}

A ferramenta executa o ciclo detectar $\rightarrow$ contextualizar
$\rightarrow$ triar $\rightarrow$ reparar $\rightarrow$ re-verificar,
ilustrado na Figura~\ref{fig:pipeline}. A análise estática (frente A) produz os
candidatos e o contexto de cada um; a etapa com modelo de linguagem (frente B)
os ordena, explica e propõe correções; a própria análise estática, com apoio de
um verificador limitado, confirma as correções.

\begin{figure}[htb]
  \centering
  \caption{Pipeline da ferramenta proposta}
  \label{fig:pipeline}
  \begin{tikzpicture}[
      node distance=5mm and 4.5mm,
      etapa/.style={draw=black, rounded corners=2pt, align=center,
        minimum height=11mm, text width=29.5mm, font=\small, fill=white},
      llm/.style={etapa, dashed},
      seta/.style={-{Stealth[length=2.2mm]}, thick, color=black},
      nota/.style={font=\scriptsize, color=black, align=center}]
    \node[etapa] (cfg) {Código C +\\configuração};
    \node[etapa, right=of cfg] (e1) {Estágio 1\\candidatos};
    \node[etapa, right=of e1] (e2) {Estágio 2\\prioridade e\\mascaramento};
    \node[etapa, right=of e2] (e3) {Estágio 3\\viabilidade (Z3)};
    \node[llm, below=12mm of e3] (tri) {Triagem\\(ordena e explica)};
    \node[llm, left=of tri] (rep) {Reparo\\com testemunha};
    \node[etapa, left=of rep] (rev) {Re-verificação};
    \node[etapa, left=of rev] (rel) {Relatório\\ordenado};
    \draw[seta] (cfg) -- (e1);
    \draw[seta] (e1) -- (e2);
    \draw[seta] (e2) -- (e3);
    \draw[seta] (e3) -- node[right, nota] {registro de\\contexto} (tri);
    \draw[seta] (tri) -- (rep);
    \draw[seta] (rep) -- (rev);
    \draw[seta] (rev) -- (rel);
    \draw[seta, dashed] (rev.north) to[bend left=18] node[left, nota, pos=0.35] {nova\\análise} (e1.south);
    \node[nota, above=1mm of e2] {só descartam com prova de impossibilidade};
    \node[nota, below=1mm of rep] {etapas com \llm\ (tracejadas): nunca descartam};
  \end{tikzpicture}
  \fonte{Elaborada pelos autores.}
\end{figure}

\section{Princípio de projeto}
\label{sec:principio}

Uma única regra orienta a ferramenta: um solver pode descartar um
candidato; o \llm, não. Um resultado \textit{UNSAT} do solver é uma prova de
que a intercalação é impossível; o julgamento de um modelo de linguagem não é.
Por isso a etapa de triagem ordena, explica e agrupa os candidatos, e o
relatório final mantém todos eles.

Essa regra não é apenas prudência. A literatura mede o efeito das duas
políticas no mesmo ponto de integração: filtros com \llm\ que descartam
candidatos perdem de 6 a 25 pontos de recall \cite{reducing-fp,adataint},
enquanto a política conservadora do LLift chega a recall de 1,00
\cite{llift}. A consequência para a avaliação é que a métrica principal não é a
precisão, mas o recall, que precisa ser de 100\%, acompanhado do
\textit{Inspection Ratio}.

\section{Entrada e configuração}
\label{sec:entrada}

A ferramenta aceita um único arquivo C ou um projeto com vários arquivos. No
segundo caso, cada arquivo é compilado separadamente para a representação
intermediária do LLVM e os resultados são ligados num só módulo com
\texttt{llvm-link}. Ligar a representação intermediária, em vez de concatenar
o código-fonte, preserva a origem de cada instrução (arquivo e linha),
resolve símbolos \texttt{static} homônimos e respeita as definições de
pré-processador de cada arquivo. A compilação usa \texttt{-g}, para manter as
informações de depuração, e \texttt{-O0}, para que nenhum acesso à memória seja
eliminado por otimização.

O modelo de interrupções é dado por um arquivo de configuração escrito
à mão, que nomeia os pontos de entrada de cada fluxo, suas prioridades e as
primitivas de mascaramento da plataforma. A alternativa de inferir esse modelo
com um \llm, como \citeonline{iris} fazem com especificações de bibliotecas, foi
considerada e rejeitada: as primitivas de mascaramento formam um conjunto
fechado por plataforma, os dois benchmarks já fornecem a configuração, e
encontrar os pontos de entrada, a parte crítica para o recall, é uma busca
textual pela API de registro de interrupções, não um julgamento. O arquivo de
configuração também limita o custo: a análise percorre apenas o código
alcançável a partir dos fluxos declarados, e não o repositório inteiro.

\section{Etapa estática}
\label{sec:etapa-estatica}

\subsection{Estágio 1: candidatos}

O primeiro estágio segue o desenho de \citeonline{intrace}, que ignora
deliberadamente condições de desvio e estado das interrupções para não perder
nada. Ele resolve quatro subproblemas:

\begin{enumerate}
  \item Identificar os fluxos. Cada ponto de entrada do arquivo de
    configuração é a raiz de uma análise de alcançabilidade. Um fluxo não
    declarado não contribui com acessos, e todo defeito que dependa dele
    desaparece sem aviso; por isso esta é a primeira premissa da ferramenta.
  \item Identificar as posições compartilhadas. Uma posição é
    compartilhada quando dois ou mais fluxos a alcançam, segundo uma análise de
    ponteiros do tipo \textit{may}. Não basta considerar variáveis globais: uma
    variável local cujo endereço é guardado em ponteiros globais também pode ser
    compartilhada.
  \item Enumerar os acessos. Para cada fluxo, percorrem-se as funções
    alcançáveis e registram-se as leituras e escritas às posições
    compartilhadas, com o fluxo, o tipo de acesso, a função, o caminho de
    chamadas, a linha do código-fonte e os laços que envolvem o acesso.
  \item Derivar os candidatos. Pares e triplas são derivados
    separadamente do mesmo conjunto de acessos, pela razão discutida na
    Seção~\ref{sec:defeitos}. As triplas se restringem aos quatro padrões não
    serializáveis da Tabela~\ref{tab:padroes}.
\end{enumerate}

Um detalhe sobre laços é decisivo para o recall. Um único acesso no código-fonte,
dentro de um laço que pode repetir, corresponde a vários acessos em execução;
por isso a mesma instrução pode ser $A_1$ e $A_2$ de uma tripla. O mesmo vale
para uma função chamada mais de uma vez pelo mesmo fluxo. O Racebench anota
defeitos exatamente com essa forma.

\subsection{Estágio 2: prioridade e mascaramento}

O segundo estágio descarta candidatos cujos fluxos comprovadamente não podem se
sobrepor. Um fluxo só é interrompido por outro de prioridade estritamente
maior, com a convenção de que número maior é prioridade maior. Fluxos de mesma
prioridade são tratados como capazes de se interromper mutuamente, pois nenhum
dos modelos formais da literatura cobre esse caso e essa é a leitura que não
perde defeitos.

O mascaramento é calculado por uma análise de fluxo de dados que determina, em
cada ponto do programa, quais interrupções estão desabilitadas. Uma interrupção
só é considerada mascarada durante o intervalo $[A_1, A_2]$ de uma tripla se
(a) estiver mascarada em todo ponto de todo caminho de $A_1$ a $A_2$ no fluxo
local, e (b) nenhum fluxo que possa executar dentro desse intervalo a reabilitar.
A segunda condição é a que evita perder o defeito do caso
\texttt{svp\_simple\_001\_001} descrito na Seção~\ref{sec:interrupcoes}. Escritas
em registradores não reconhecidas nunca contam como mascaramento.

Nenhum modelo de periodicidade é assumido: o Racebench declara que uma
interrupção pode ocorrer em qualquer ponto habilitado, qualquer número de
vezes. Para pares, a pergunta é se $B$ pode executar no \textit{instante} de
$A$; para triplas, se pode executar em algum ponto do \textit{intervalo}
$[A_1, A_2]$. As duas perguntas são respondidas separadamente.

\subsection{Estágio 3: viabilidade com Z3}

O terceiro estágio aplica a análise de viabilidade de caminho de
\citeonline{intrace} com o solver Z3. Como o estágio 2, só descarta um candidato com prova, isto é, com resultado \textit{UNSAT}.
O resultado tem quatro valores (\textit{UNSAT}, \textit{SAT}, \textit{inconclusivo} e \textit{não executado}), e um resultado inconclusivo
(por limite de tempo ou construção não suportada) nunca é tratado como decidido:
o candidato segue para a triagem, com a informação de até onde o solver chegou.

A posição desse estágio antes do \llm\ é fixa por três razões: é um filtro correto, que só descarta com prova; reduz fortemente o volume (no IntRace, cerca de 208
candidatos por programa caem para cerca de 95 após o segundo estágio, e a
eliminação de 86,2\% relatada para o terceiro deixaria cerca de 13 \cite{intrace}), o que torna a triagem por \llm\ economicamente viável; e reproduz a arquitetura do LLift, em que o \llm\ recebe
apenas os casos que a execução simbólica não conseguiu decidir \cite{llift}.

\section{O registro de contexto}
\label{sec:registro}

Para cada candidato que sobrevive à etapa estática, a ferramenta emite um
registro de contexto: um documento JSON com as evidências necessárias
para julgá-lo. \citeonline{llift} atribuem 5 de seus 13 falsos positivos a
lacunas no relatório entregue pelo analisador ao modelo, o que justifica tratar
esse registro como um produto de primeira classe da análise, e não como um
subproduto. A Tabela~\ref{tab:registro} resume seus campos.

\begin{table}[htb]
  \centering
  \caption{Campos do registro de contexto}
  \label{tab:registro}
  \begin{tabular}{@{}lL{11cm}@{}}
    \toprule
    Campo & Conteúdo \\
    \midrule
    classe & par (\textit{data race}) ou tripla (violação de atomicidade) \\
    variável & nome, tipo, se é \texttt{volatile}, local da declaração \\
    acessos & para cada um: leitura ou escrita, linha, função e fluxo \\
    fluxos & ponto de entrada, prioridade e como o fluxo foi identificado \\
    caminhos & sequência de chamadas do ponto de entrada até cada acesso \\
    mascaramento & estado de cada interrupção em cada acesso e no intervalo \\
    laços & laços que envolvem cada acesso; se $A_1$ e $A_2$ são a mesma instrução \\
    código & corpo completo das funções que contêm os acessos \\
    solver & veredicto do Z3 e, se inconclusivo, o motivo \\
    proveniência & quais fatos foram provados, quais vêm da configuração, quais são desconhecidos \\
    \bottomrule
  \end{tabular}
  \fonte{Elaborada pelos autores.}
\end{table}

O registro não traz um fatiamento de programa pré-calculado. Um
fatiamento sobre uma variável global tende a abarcar quase todo o programa, e um
fatiamento incompleto leva o modelo a declarar, com confiança, que um defeito
real é impossível. Em vez disso, o registro traz as camadas sempre úteis, e o
modelo pede o que falta. O campo de proveniência serve ao mesmo propósito:
evidência parcial é marcada como parcial.

\section{Etapa com modelo de linguagem}
\label{sec:etapa-llm}

\subsection{Triagem: duas perguntas, quatro grupos}

O modelo nunca escolhe diretamente onde o candidato fica. Ele responde a duas
perguntas estreitas: a intercalação pode acontecer? (viabilidade) e,
se puder, ela causa dano? (nocividade). Separar as perguntas evita que
``benigno'' se transforme silenciosamente em ``impossível''. O grupo do
candidato é derivado das duas respostas por uma tabela fixa
(Tabela~\ref{tab:grupos}), na qual qualquer dúvida sobe o candidato na fila.

\begin{table}[htb]
  \centering
  \caption{Derivação do grupo de cada candidato, em ordem de leitura}
  \label{tab:grupos}
  \begin{tabular}{@{}cL{3.6cm}L{6.0cm}L{3.4cm}@{}}
    \toprule
    Ordem & Grupo & Viabilidade & Nocividade \\
    \midrule
    1 & provavelmente real & possível & nociva \\
    2 & incerto & incerta, ou possível com nocividade incerta & qualquer \\
    3 & provavelmente benigno & possível & benigna \\
    4 & provavelmente inviável & impossível, com o elemento que impede & não perguntada \\
    \bottomrule
  \end{tabular}
  \fonte{Elaborada pelos autores.}
\end{table}

Todos os grupos aparecem no relatório. O grupo só determina a posição de leitura,
e é essa ordem que o \textit{Inspection Ratio} avalia.

\subsection{Desenho do \textit{prompt}}

O desenho segue os quatro componentes de \citeonline{llift}, adaptados ao
domínio:

\begin{itemize}
  \item Regras do domínio por exemplos. O conhecimento prévio do modelo
    vem sobretudo de \textit{threads} e o leva a procurar travas e relações
    \textit{happens-before}. O \textit{prompt} ensina uma tabela de padrões de
    interrupção com seus veredictos: uma seção crítica cobrindo todo o intervalo
    torna a tripla inviável; duas seções separadas deixam uma lacuna em que a
    \isr\ cabe; prioridade estritamente menor impede a preempção, e prioridade
    igual não; a mesma instrução num laço pode formar uma tripla consigo mesma.
  \item Pedidos progressivos de contexto. O modelo pode pedir, num
    formato fixo, a definição de uma função, o corpo de uma \isr, o estado de
    mascaramento numa linha ou a expansão de uma macro. O registro dos pedidos
    mede, na prática, que contexto o registro deveria trazer.
  \item Decomposição da tarefa. Viabilidade e nocividade são perguntadas
    em conversas separadas.
  \item Autovalidação enviesada para o recall. Antes de responder, o
    modelo confere regras como: mascaramento desconhecido conta como
    habilitado; uma função cuja definição não foi obtida pode acessar a
    variável; ``incerto'' é sempre uma resposta segura.
\end{itemize}

O modelo raciocina antes de dar o veredicto, e a resposta final é uma saída
estruturada com a explicação antes do veredicto, como recomendam
\citeonline{iris}. O código analisado é delimitado e declarado como
dado, nunca instrução: um comentário no código que tente convencer o
modelo a ignorar um defeito não pode ter efeito \cite{sastgenius}.

A etapa é construída de forma independente do provedor do modelo. Isso permite
manter o \textit{prompt} fixo e trocar o modelo, e assim verificar neste domínio a
afirmação de \citeonline{llift} de que a arquitetura do \textit{prompt} pesa mais
que a escolha do modelo.

\subsection{Reparo com testemunha}

Para cada candidato, o modelo pode propor uma correção no vocabulário
específico de interrupções enumerado por \citeonline{sdracer}: desabilitar e
reabilitar a interrupção em volta do acesso, estender uma seção crítica
existente ou fundir seções adjacentes. O escopo depende da classe do defeito:
um \textit{data race} se corrige protegendo um acesso; uma violação de
atomicidade, tornando todo o intervalo $[A_1, A_2]$ atômico.

Inspirada na validação por \textit{exploits} de \citeonline{sastgenius}, toda
proposta de correção exige uma testemunha: qual \isr\ dispara, em que
ponto, em que ordem ocorrem os acessos e que estado resulta diferente de
qualquer execução em série. A testemunha é conferida mecanicamente contra o
registro de contexto antes de a correção ser aceita.

\section{Re-verificação}
\label{sec:reverificacao}

Uma correção é verificada em duas etapas, que falham em direções opostas.
Primeiro, a própria análise estática é executada de novo sobre o programa
corrigido. Como seus filtros só descartam com prova, o desaparecimento do
candidato é uma certificação, sob as premissas declaradas, de que o defeito
sumiu; e a nova análise completa é a única forma de detectar
regressões, como uma seção crítica estendida que cria outro defeito.
Se o candidato persistir, o resultado é inconclusivo, pois uma análise
superaproximada pode continuar reportando um defeito já corrigido. Segundo, os
casos inconclusivos são resolvidos por um verificador limitado (o CBMC \cite{cbmc} ou o próprio BMC4AV, cujo código-fonte está disponível), codificando a violação como uma asserção. Por fim, como observado por
\citeonline{sdracer}, verifica-se se a correção alterou o fluxo de controle ou
a latência, que em código de interrupção é uma propriedade de correção.

\section{Interface}
\label{sec:interface}

A ferramenta é pensada para ser operada por um painel, e não apenas por linha
de comando: configurar um programa, executar a análise, ler cada candidato ao
lado do código-fonte e aceitar ou rejeitar uma correção. A literatura compete,
na prática, em automação: cada ferramenta nova ataca o trabalho manual exigido pela anterior \cite{intrace,nichecker,bmc4av}. Por isso, a usabilidade é tratada aqui como parte da contribuição.

\section{Contratos entre as etapas}
\label{sec:contratos}

Para que as duas frentes possam avançar em paralelo, as interfaces entre elas
são definidas como esquemas JSON, com exemplos (Tabela~\ref{tab:contratos}).
\ifparcial\else O Apêndice~\ref{ap:contratos} os detalha.\fi

\begin{table}[htb]
  \centering
  \caption{Contratos entre as partes da ferramenta}
  \label{tab:contratos}
  \begin{tabular}{@{}llll@{}}
    \toprule
    & Contrato & Produz & Consome \\
    \midrule
    C1 & arquivo de configuração (modelo de interrupções) & frente A & frentes A e C \\
    C2 & registro de contexto de cada candidato & frente A & frentes B e C \\
    C3 & organização de uma execução em disco & frente A & frentes B e C \\
    C4 & pedidos de contexto e respostas & frente B & frente A \\
    \bottomrule
  \end{tabular}
  \fonte{Elaborada pelos autores.}
\end{table}

\section{Premissas e ameaças ao recall}
\label{sec:premissas}

A afirmação de que a ferramenta não perde defeitos vale em relação a uma lista
explícita de premissas, mantida junto com o código. As principais são:

\begin{itemize}
  \item o arquivo de configuração nomeia todos os fluxos; um tratador não
    declarado é invisível, e esta é a maior ameaça ao recall;
  \item uma chamada indireta cujos alvos não se resolvem pode chamar qualquer
    função de assinatura compatível cujo endereço foi tomado;
  \item posições compartilhadas são identificadas por análise \textit{may}, e a
    correção da análise de ponteiros do SVF é herdada, não reverificada;
  \item uma função sem corpo no programa pode ler e escrever qualquer variável
    global;
  \item fluxos de mesma prioridade podem se interromper, e o aninhamento é
    permitido;
  \item uma interrupção só está mascarada se estiver mascarada em todos os
    caminhos, e nenhum outro fluxo a reabilitar no intervalo;
  \item só uma prova descarta um candidato; o resultado inconclusivo do solver e
    o julgamento do \llm\ nunca descartam.
\end{itemize}

Uma premissa é deliberadamente não conservadora: por padrão, uma \isr\
não interrompe a si mesma. O Racebench não exclui a reentrância e a literatura
diverge sobre ela; a escolha fica registrada para revisão antes de qualquer
afirmação final de recall.
